% Chapter 3 -- The Copy Command (AN-08 opener, AN-09 entry; annex Sec 7).
% Sources (Law 1): os/i123/app/face_a.c menu_select / range_accept (the
% FROM and TO prompts, the single-cell TO rule and its message),
% os/i123/core/range.c i123_transfer (what is copied, what the destination
% loses, reference words kept: relative words are offsets, absolute words
% coordinates), spec/i123/era_shortfalls.md I123-P3-07 (TO ranges larger
% than one cell refuse).  Menu tree generated from the shipped resource.
% GATES: test-i123-copy-move-model test-i123-menu-tree test-pubs-i123 test-pubs-screens
\Gates{test-i123-copy-move-model, test-i123-menu-tree, test-pubs-i123, test-pubs-screens}
\chapter[The Copy Command]{\textit{The Copy Command}}
\input{i123-tree-Copy}

The Copy command copies the contents of a cell or a range of cells to
another part of the worksheet, together with each cell's format and
protection setting.

\begin{LotTable}{p{30mm}p{98mm}}
\LotHead{\textbf{Command} & \textbf{Task}}
/Copy & Copies a cell or a range to another location in the worksheet\LotEnd
\end{LotTable}

\section{Common Uses of the Copy Command}
\begin{itemize}
\item Entering the same formula in a row or a column of cells
  (\slashcmd{/Copy})
\item Repeating a heading, a label or a block of numbers elsewhere in the
  worksheet (\slashcmd{/Copy})
\item Making a working copy of a table before you change it
  (\slashcmd{/Copy})
\end{itemize}

\section{Related Reading}
\begin{itemize}
\item ``Formulas'' and ``Ranges'' in Chapter~1
\item Chapter~7, ``The Move Command''
\item \slashcmd{/Range Value} and \slashcmd{/Range Trans} in Chapter~10,
  which copy the current values of formulas instead of the formulas
\item Lesson~5 in the \doctitle{Tutorial}
\end{itemize}

\section{Terms You Need to Know}
\begin{itemize}
\item The \term{FROM range} is the cell or range you copy.
\item The \term{TO range} is where the copy goes.  In Initech 123 you name
  the TO range by its upper left cell.
\item A \term{relative reference} in a formula, such as \texttt{A1}, refers
  to a cell by its position relative to the formula.  When you copy the
  formula, the reference refers to the cell in the same position relative
  to the copy.
\item An \term{absolute reference}, such as \texttt{\$A\$1}, always refers
  to the same cell, wherever the formula is copied.
\end{itemize}

\section{Before You Use /Copy}
\begin{Caution}
/Copy replaces whatever the TO range contains.  A cell in the TO range that
is not part of the copy keeps nothing: the copy of a blank cell is a blank
cell.  Save the worksheet before you copy over data you may want back.
\end{Caution}

\MenuEntry{/Copy}
/Copy copies a cell or a range of cells, with their formats and protection
settings, to another part of the worksheet.

\FlairShot[fig:copy-to]{copyto}{The TO prompt of /Copy, with the FROM range
B1..C1 already given}

\BHead{Terms You Need to Know}
\begin{itemize}
\item The \term{FROM range} is the cell or range you copy.
\item The \term{TO range} is the place the copy goes, named by its upper
  left cell.
\end{itemize}
Initech 123 copies every cell of the FROM range to the cell in the same
position relative to the upper left cell of the TO range.  The rectangle
the copy covers is first cleared: whatever was there, including a cell's
format, is replaced by the copy.  A label or a number is copied as it is.
A formula is copied with its references: a relative reference keeps its
distance from the formula, so a formula that adds the two cells to its left
still adds the two cells to its left in the copy; an absolute reference
keeps naming the same cell.

When the worksheet is set to recalculate automatically (the setting
\slashcmd{/Worksheet Global Recalculation Automatic}, which Initech 123 uses
until you change it), the copied formulas are calculated as soon as the
copy is made.

\BHead{Procedure}
\begin{Steps123}
\item Select /Copy.  The control panel shows
  \texttt{Enter range to copy FROM:} followed by the address of the cell
  pointer.
\item Specify the FROM range and press \kname{ENTER}.

  Type the range, for example \type{B1..C1}, or a range name; or press
  \kname{ENTER} at once to copy just the cell at the cell pointer.  You can
  also move the cell pointer with the pointer-movement keys to the one cell
  you want to copy, and press \kname{ENTER}.
\item Specify the TO range and press \kname{ENTER}.

  Type the address of the upper left cell of the place the copy goes, or
  move the cell pointer there, and press \kname{ENTER}.  Initech 123 makes
  the copy and returns to \ind{READY} mode.
\end{Steps123}

\FlairShot[fig:copy-done]{copydone}{The formulas in B1..C1 copied to D1..E1}

Figure~\ref{fig:copy-done} shows the result of copying B1..C1, which hold
the formulas \texttt{-A1*2} and \texttt{-B1*2}, to D1.  The copy in D1
refers to C1, the cell to its left, and the copy in E1 to D1, just as the
originals refer to the cells to their left.

\LabelHead{Tip}
\begin{itemize}
\item To copy a formula to several places, copy it once to each place,
  naming each upper left cell in turn.
\item To keep a reference fixed when you copy a formula, make it absolute
  with dollar signs, as in \texttt{\$B\$1}, before you copy.
\end{itemize}

\begin{Note}
Give the TO range as a single cell.  If you type a TO range of more than one
cell, such as \type{D1..E3}, Initech 123 displays
\texttt{Specify the upper left destination cell} and makes no copy; press
\kname{ESC} to return to the worksheet.  The size of the copy is always the
size of the FROM range.
\end{Note}

\begin{Note}
If the worksheet is protected (\slashcmd{/Worksheet Global Protection
Enable}) and the TO range includes a protected cell, Initech 123 displays
\texttt{Command refused} and makes no copy.
\end{Note}

\begin{Note}
The FROM and TO ranges may overlap.  Initech 123 takes the contents of the
FROM range as they were before the copy began, so every cell of the copy
receives the original contents.
\end{Note}

\section{Copying a Block with Mixed References: An Example}
A \term{mixed reference} fixes only the column (\texttt{\$A1}) or only the
row (\texttt{A\$1}).  B1 holds \texttt{+A1*\$A\$1} and C1 holds
\texttt{+A\$1*\$A1}; B1..C1 is copied to B2.  The copy of B1 keeps the
absolute \texttt{\$A\$1} and moves \texttt{A1} down a row; the copy of C1
keeps row 1 in \texttt{A\$1} and column A in \texttt{\$A1}, whose row
moves down, so both copies in row 2 multiply A2 by A1:

\LotExample{examples/copymixed}

\section{Copying a Formula Down a Column: An Example}
In this example a column of prices in B2..B4 and a tax rate in B6 are used
to compute the tax on each price.  The tax rate is referred to absolutely,
so that every copy of the formula uses it.
\begin{Steps123}
\item Enter the prices: \type{100} in B2, \type{250} in B3 and \type{40} in
  B4; enter \type{.08} in B6.
\item Move the cell pointer to C2 and type \type{+B2*\$B\$6}, then press
  \kname{ENTER}.  C2 shows \texttt{8}.
\item Select /Copy, press \kname{ENTER} to copy C2, type \type{C3} and press
  \kname{ENTER}.
\item The cell pointer is still on C2.  Select /Copy again, press
  \kname{ENTER}, type \type{C4} and press \kname{ENTER}.
\end{Steps123}
The entries and the three formula cells, as the control panel and the
worksheet show them afterwards:

\LotExample{examples/copytax}
